\documentclass[11pt]{article}
\usepackage[T1]{fontenc}
\usepackage[utf8]{inputenc}
\usepackage{lmodern}
\usepackage{amsmath,amssymb,amsthm,mathtools}
\usepackage{booktabs,array}
\usepackage[margin=1in]{geometry}
\usepackage{microtype}
\usepackage{xcolor}
\usepackage{hyperref}
\usepackage{xurl}
\hypersetup{
  colorlinks=true,
  linkcolor=blue!55!black,
  citecolor=blue!55!black,
  urlcolor=blue!55!black,
  pdftitle={Proper transposed sesqui arrays at all Sylvester--Hadamard powers},
  pdfauthor={Carptopus}
}
\newtheorem{theorem}{Theorem}[section]
\newtheorem{lemma}[theorem]{Lemma}
\newtheorem{proposition}[theorem]{Proposition}
\newtheorem{corollary}[theorem]{Corollary}
\theoremstyle{definition}\newtheorem{definition}[theorem]{Definition}
\theoremstyle{remark}\newtheorem{remark}[theorem]{Remark}
\DeclareMathOperator{\Tr}{Tr}
\DeclareMathOperator{\GF}{GF}
\newcommand{\F}{\mathbb F}
\newcommand{\SA}{SA^{\mathsf T}}
\newcommand{\ind}{\mathbf 1}
\title{Proper transposed sesqui arrays at all Sylvester--Hadamard powers}
\author{Carptopus\\\href{mailto:carptopus@163.com}{\texttt{carptopus@163.com}}}
\date{18 August 2026}

\begin{document}
\maketitle

\begin{abstract}
For every power of two \(t=2^k\), \(k\geq2\), we construct a proper
transposed sesqui array
\[
 \SA(4t-2,t,-,t-1,t:(2t-1)\times2t).
\]
The columns form the classical Sylvester--Hadamard trace design.  We reduce
the cell-ordering problem to a two-to-one finite-field map.  Odd field
dimensions are handled by an explicit inverse-pair map.  In even dimension a
Subiaco map and its zero fibre reduce the problem to two fixed algebraic
curves.  Geometric irreducibility, Artin--Schreier trace covers, and a
Hasse--Weil estimate cover all even dimensions at least \(34\); fifteen exact
polynomial-gcd certificates close the remaining even dimensions.  A fixed
resultant certificate proves properness, and an explicit array closes \(t=4\).
\end{abstract}

\paragraph{Status.}
The proof and its exact computations have passed implementation-separated
internal checks.  External mathematical review and an author-level priority
inquiry remain pending, so all priority claims are qualified.

\section{Introduction}
A binary equireplicate row--column array is a transposed sesqui array when
column--column and row--column concurrences are constant.  It is proper when
row--row concurrence is nonconstant.  We study
\begin{equation}
 (v,e,\lambda_{CC},\lambda_{RC};r,c)
 =(2q-2,q/2,q/2-1,q/2;q-1,q),\qquad q=2^d,
 \label{eq:qparams}
\end{equation}
or, with \(q=2t\),
\begin{equation}
 \SA(4t-2,t,-,t-1,t:(2t-1)\times2t).
 \label{eq:tparams}
\end{equation}

The columns below are classical affine trace hyperplanes.  In the Sylvester
case they agree block by block with the column incidence in the resolvable
unordered triple-array construction of Gordeev and Öhman
\cite[Theorem~5.13]{gordeev-ohman}.  That prior construction has constant
row concurrence.  The issue addressed here is the compatible ordered
\emph{proper} array.

The proof has five parts: a general two-fibre ordering theorem; an explicit
odd-dimensional family; a Subiaco zero-fibre bridge in even dimension; a
fixed-curve Hasse--Weil argument for all large even dimensions; and exact
certificates for the finite complement.  The computation is fixed and finite:
two resultants, five irreducibility and trace checks, fifteen finite-field
witnesses, and one boundary array.  It contains no array classification or
extrapolation from sampled dimensions.

To the best of our knowledge, \eqref{eq:tparams} is the first proper
construction for every power of two \(t\).  The \(t=8\) instance fills the
proper entry marked \(?,T\) by Jäger et al.~\cite[Table~2]{jager}; \(T\)
records a triple array rather than a proper one.

\section{Trace columns and a two-fibre ordering theorem}
Let \(F=\GF(q)\), \(q=2^d\), and let \(\Tr:F\to\F_2\) be the absolute
trace.  Index symbols by \((h,b)\in F^*\times\F_2\).  For \(c\in F\), set
\begin{equation}
 \mathcal B_c=\{(h,\Tr(ch)):h\in F^*\}.
 \label{eq:columns}
\end{equation}

\begin{lemma}
For \(c\neq c'\), \(|\mathcal B_c\cap\mathcal B_{c'}|=q/2-1\).
\end{lemma}
\begin{proof}
A common symbol satisfies \(\Tr((c+c')h)=0\).  The nondegenerate trace
pairing gives a hyperplane of size \(q/2\); remove \(h=0\).
\end{proof}

\begin{theorem}[Two-fibre ordering]
\label{thm:twofibre}
Suppose \(f:F\to F^*\) is exactly two-to-one with image \(Y\), the map
\(x\mapsto xf(x)\) permutes \(F^*\), and every fibre \(\{x,x'\}\) satisfies
\begin{equation}
 \Tr((x+x')f(x))=1.
 \label{eq:fibretrace}
\end{equation}
For rows \(r\in F^*\), columns \(c\in F\), define
\begin{equation}
 A(r,c)=\bigl(rf(cr),\Tr(crf(cr))\bigr).
 \label{eq:cells}
\end{equation}
Then \(A\) is a binary equireplicate transposed sesqui array with
\eqref{eq:qparams}, and for distinct rows
\begin{equation}
 |A_r\cap A_s|=2|Y\cap(s/r)Y|.
 \label{eq:rr}
\end{equation}
\end{theorem}
\begin{proof}
For \(c=0\), the first coordinate \(rf(0)\) runs through \(F^*\).  For
\(c\neq0\), putting \(x=cr\) gives first coordinate \(xf(x)/c\), which also
runs through \(F^*\).  Thus column \(c\) is \(\mathcal B_c\).

In a fixed row, the two inputs in a fibre of value \(y\) give first coordinate
\(ry\).  Condition \eqref{eq:fibretrace} gives opposite binary coordinates,
so the row contains both symbols over every member of \(rY\), with no
repetition.  A fixed first coordinate \(h\) is supported by the \(q/2\) rows
\(r=h/y\), \(y\in Y\); hence each symbol has replication \(q/2\).  A row
and column share one symbol over each member of \(rY\), so their intersection
is \(q/2\).  Two rows share both symbols exactly over \(rY\cap sY\), proving
\eqref{eq:rr}.
\end{proof}

\section{Odd field dimension}
Assume \(d\) is odd and define
\begin{equation}
 f(0)=f(1)=1,\qquad
 f(x)=\frac1{x+x^{-1}}=\frac{x}{x^2+1}
 \quad(x\notin\{0,1\}).
 \label{eq:inverse}
\end{equation}
If \(H=\ker\Tr\), then
\begin{equation}
 Y=(H\setminus\{0\})\cup\{1\}.
 \label{eq:Yodd}
\end{equation}
Indeed, \(f(x)=y\) reduces through \(u=xy\) to
\(u^2+u=y^2\).  It has two roots precisely for nonzero \(y\in H\), and the
roots in \(x\) are mutual inverses.  The special fibre \(\{0,1\}\) has image
one and is disjoint from the ordinary images because \(\Tr(1)=1\).  Every
fibre satisfies \eqref{eq:fibretrace}.  Moreover,
\[
 xf(x)=\left(\frac{x}{x+1}\right)^2
\]
on \(F^*\setminus\{1\}\), so the required permutation follows from a Möbius
map and Frobenius.

For \(a\in F^*\setminus\{1\}\),
\begin{equation}
 |Y\cap aY|
 =q/4-1+\ind_{\Tr(a)=0}+\ind_{\Tr(a^{-1})=0}.
 \label{eq:oddintersection}
\end{equation}
Writing
\[
 K_q=\sum_{a\in F^*}(-1)^{\Tr(a+a^{-1})},
\]
elementary character sums give
\[
 N_{00}=(q-3+K_q)/4,\qquad N_{11}=(q+1+K_q)/4.
\]
The binary Kloosterman bound \(|K_q|\leq2\sqrt q\)
\cite{conrad-kloosterman} gives \(N_{00}>0\) and
\(N_{11}>1\) for \(q\geq32\).  Thus
\eqref{eq:oddintersection} assumes at least two values.

\begin{theorem}
\label{thm:odd}
For every odd \(d\geq5\), \eqref{eq:cells} and \eqref{eq:inverse} give a
proper array with \eqref{eq:qparams}.
\end{theorem}

\section{The Subiaco bridge in even dimension}
Let \(d\geq4\) be even.  Call \(c\in F\) admissible when
\begin{equation}
 c\neq0,\qquad c^2+c+1\neq0,\qquad \Tr(c^{-1})=1.
 \label{eq:admissible}
\end{equation}
Put
\[
 u=1+c+c^2,\qquad D(x)=x^2+cx+1,
\]
\[
 N_F(x)=c^2(x^4+x)+c^2u(x^3+x^2),
\]
and take the normalized Subiaco o-polynomial
\begin{equation}
 F_c(x)=\frac{N_F(x)}{D(x)^2}+x^{1/2},
 \qquad B_c(x)=F_c(x)+ux.
 \label{eq:subiaco}
\end{equation}
The denominator has no field root under \eqref{eq:admissible}: if \(D(x)=0\),
then \(y=x/c\) would satisfy \(y^2+y=c^{-2}\), whereas
\(\Tr(c^{-2})=\Tr(c^{-1})=1\).  We make the finite-geometric input and the
additional algebra explicit.  Put
\[
 \eta=c^2+c^5+c^{1/2},\qquad \kappa=\frac{\eta}{cu},
\]
Here \(u\ne0\), and also \(\eta\ne0\).  For the latter claim, suppose
\(\eta=0\) and put \(w=c^{-1/2}\).  Then
\[
 w^9+w^6+1=(w^3+w+1)(w^6+w^4+w^2+w+1)=0.
\]
The two displayed factors are irreducible over \(\mathbb F_2\), of degrees
three and six.  In the cubic case, \(3\mid d\); since \(d\) is even, the
extension multiplier \(d/3\) is even, so \(\Tr_F(w^2)=0\).  In the sextic
case, \(6\mid d\), and the zero \(w^5\)-coefficient gives
\(\Tr_{\mathbb F_{64}/\mathbb F_2}(w)=0\), hence again
\(\Tr_F(w^2)=0\).  Both cases contradict
\(\Tr(c^{-1})=\Tr(w^2)=1\).

The exact coordinate formulas below are recorded by O'Keefe--Thas
\cite[pp.~178--179]{okeefe-thas}, with their \(d\) equal to our \(c\) and their scalar
\(a\) equal to our \(\kappa\):
\[
 N_G(x)=c^4x^4+c^3(1+c^2+c^4)x^3+c^3(1+c^2)x,
\]
and
\begin{equation}
 G_c(x)=\frac{N_G(x)}{\eta D(x)^2}
       +\frac{c^{1/2}x^{1/2}}{\eta}.
 \label{eq:subiacocompanion}
\end{equation}
The Subiaco construction \cite{flocks-ovals} supplies that \(F_c\) is an
o-polynomial and that, in the normalized \(q\)-clan convention
\cite[Eq.~(19)]{bader-qclan}, the matrices
\[
 A_x=\begin{pmatrix}F_c(x)&x^{1/2}\\0&\kappa G_c(x)\end{pmatrix}
 \qquad(x\in F)
\]
form a \(q\)-clan.  Thus for \(s\ne t\) the binary quadratic form attached
to \(A_s+A_t\) is anisotropic.  Its cross coefficient has square \(s+t\),
so the characteristic-two anisotropy criterion gives exactly
\begin{equation}
 \Tr\left(
   \kappa\frac{(F_c(s)+F_c(t))(G_c(s)+G_c(t))}{s+t}
 \right)=1.
 \label{eq:qclan}
\end{equation}
Here \(F_c(0)=0\) and \(F_c(1)=1\), so the first fact and the
o-polynomial/two-to-one criterion \cite[Theorem~2.2]{kolsch} imply that
\(B_c(x)=F_c(x)+ux\) is exactly two-to-one.

\begin{lemma}[Subiaco fibres]
\label{lem:subfibres}
Every fibre \(\{s,t\}\) of \(B_c\) satisfies
\begin{equation}
 \Tr((s+t)/c)=1.
 \label{eq:subtrace}
\end{equation}
\end{lemma}
\begin{proof}
Define
\[
 P(x)=c^2x(x+c),\qquad
 H(x)=\frac{P(x)}{D(x)}+(u+c^{-1/2})x^{1/2}.
\]
Direct expansion in characteristic two gives
\[
 N_G(x)/c+uN_F(x)=P(x)^2+P(x)D(x),
\]
and hence
\begin{equation}
 \kappa uG_c(x)+x/c+uF_c(x)+u^2x=H(x)^2+H(x).
 \label{eq:artinschreierbridge}
\end{equation}
Taking absolute traces in \eqref{eq:artinschreierbridge} yields
\[
 \Tr(\kappa uG_c(x)+x/c)=\Tr(uB_c(x)).
\]
If \(B_c(s)=B_c(t)\), addition at \(s,t\) gives
\[
 \Tr\bigl(\kappa u(G_c(s)+G_c(t))+(s+t)/c\bigr)=0.
\]
But \(F_c(s)+F_c(t)=u(s+t)\), so \eqref{eq:qclan} gives
\(\Tr(\kappa u(G_c(s)+G_c(t)))=1\).  The last two identities prove
\eqref{eq:subtrace}.
\end{proof}

\begin{lemma}[Exterior-line bridge]
\label{lem:bridge}
Let \(I=B_c(F)\).  For a fibre \(\{a,a'\}\), put
\(b=B_c(a)\), \(\delta=a+a'\).  If
\begin{equation}
 u\delta\notin I+b,
 \label{eq:exterior}
\end{equation}
then Theorem~\ref{thm:twofibre} applies to
\begin{equation}
 x_z=\frac{z+a}{B_c(z)+b},\qquad
 f(x_z)=c^{-1}(B_c(z)+b)
 \label{eq:change}
\end{equation}
for \(z\notin\{a,a'\}\), completed by
\(f(0)=f(u^{-1})=c^{-1}u\delta\).
\end{lemma}
\begin{proof}
Set
\[
 r_a(z)=\frac{F_c(z)+F_c(a)}{z+a},\qquad r_a(a)=0.
\]
The secant-slope property of an o-polynomial says that \(r_a\) permutes
\(F\), and
\[
 B_c(z)+b=(z+a)(r_a(z)+u).
\]
The partner \(a'\) is the unique point with \(r_a(a')=u\).  Hence the
\(x_z\) in \eqref{eq:change} run through
\(F\setminus\{0,u^{-1}\}\).  Every nonzero member of \(I+b\) gives one
ordinary two-element label fibre, while \eqref{eq:exterior} prevents the
special label from colliding with an ordinary one.  Thus \(f\) is exactly
two-to-one.

For ordinary inputs,
\[
 x_zf(x_z)=c^{-1}(z+a).
\]
These products cover \(F\setminus\{0,c^{-1}\delta\}\); the two special
inputs give the missing products \(0,c^{-1}\delta\).  Thus \(xf(x)\)
permutes \(F\).  Finally, an ordinary fibre coming from
\(B_c(s)=B_c(t)\) satisfies \eqref{eq:fibretrace} by
Lemma~\ref{lem:subfibres}, while the special fibre satisfies it because
\(\Tr(\delta/c)=1\) for the anchor pair \(a,a'\).
\end{proof}

\section{The zero fibre and a fixed properness certificate}
Write \(c=h^2\).  The fibre of \(B_c\) containing zero is
\begin{equation}
 \{0,cr^2\},\qquad \Tr(r)=1,
 \label{eq:zerofibre}
\end{equation}
where the nonzero partner satisfies
\begin{align}
K_h(R)={}&1+(h+h^3)R+(h^7+h^9+h^{11})R^3+h^8R^4\notag\\
&+h^{11}R^7+h^8R^8+(h^9+h^{11}+h^{13})R^9=0.
\label{eq:Kh}
\end{align}
The exterior condition becomes
\begin{equation}
 ucr^2\notin I.
 \label{eq:zeroexterior}
\end{equation}

For \(Y\subset F^*\), put \(S_j(Y)=\sum_{y\in Y}y^j\).  If
\(|Y|=q/2\) and every nontrivial \(|Y\cap aY|\) equals \(q/4\), then \(Y\)
is a cyclic \((v,k,\lambda)=(q-1,q/2,q/4)\) difference set.  Its order is
\(n=k-\lambda=q/4\), a power of the prime \(2\), and
\(\gcd(2,v)=\gcd(2,q-1)=1\).  The prime-power case of the Second Multiplier
Theorem \cite[Section~4]{xiang} therefore makes \(2\) a numerical multiplier,
so \(Y^2=\alpha Y\) for some \(\alpha\in F^*\).
When \(S_1(Y)\neq0\), summing and normalizing gives
\begin{equation}
 S_j(Y)\in\{0,S_1(Y)^j\}\qquad\text{for all }j.
 \label{eq:powerobstruction}
\end{equation}

Direct trace-hyperplane summation gives, for \(I=B_c(F)\),
\begin{align}
 A:=\sum_{v\in I}v&=cu,\label{eq:S1}\\
 C:=\sum_{v\in I}v^3
 &=h^3+h^7+h^9+h^{10}+h^{11}+h^{12}+h^{16}\notag\\
 &=h^3(h^2+h+1)(h^3+h+1)^2(h^5+h^4+h^2+h+1).
 \label{eq:S3}
\end{align}
Here every factor is nonzero under \eqref{eq:admissible}.  Indeed \(h\ne0\),
\((h^2+h+1)^2=c^2+c+1=u\ne0\), and the remaining cubic and quintic factors
are irreducible over \(\F_2\).  If either vanished in \(F=\F_{2^d}\), then its
degree \(m\in\{3,5\}\) would divide \(d\).  Since \(d\) is even, \(d/m\) would
be even, and every element of \(\F_{2^m}\) would have absolute trace zero in
\(F\); applied to \(h^{-2}=c^{-1}\), this contradicts
\eqref{eq:admissible}.  For the zero fibre the unscaled label set is
\[
 X=(I\setminus\{0\})\cup\{ucr^2\}.
\]
Since \(d\) is even and \(\Tr(r)=1\), one has \(r\ne1\); hence
\(S_1(X)=cu(1+r^2)\ne0\), as required when applying
\eqref{eq:powerobstruction} below.

\begin{lemma}[Computer-assisted fixed resultant]
\label{lem:resultant}
For every even \(d\geq4\) and admissible \(c=h^2\), the zero fibre cannot
satisfy \(S_3(X)=0\) or \(S_3(X)=S_1(X)^3\).
\end{lemma}
\begin{proof}
Eliminating \(R\) between \eqref{eq:Kh} and the two degeneracy equations gives
\begin{equation}
 h^{54}(h+1)^2(h^2+h+1)^{14}R_{21}(h)R_{109}(h)
 \label{eq:res0}
\end{equation}
and
\begin{equation}
 h^{54}(h+1)^2(h^2+h+1)^{14}
 R_{26}(h)R_{29}(h)R_{75}(h).
 \label{eq:res1}
\end{equation}
The five \(R_m\) are monic irreducible polynomials over \(\F_2\).  Their
degrees and inverse-square root traces are
\begin{center}
\begin{tabular}{c@{\quad}ccccc}
\toprule
factor & \(R_{21}\) & \(R_{109}\) & \(R_{26}\) & \(R_{29}\) & \(R_{75}\)\\
\midrule
degree & 21 & 109 & 26 & 29 & 75\\
\(\Tr(\alpha^{-2})\) & 0 & 0 & 0 & 1 & 0\\
\bottomrule
\end{tabular}
\end{center}
The low-degree factors contradict admissibility.  Each trace-zero factor does
so in every ambient extension.  If \(R_{29}(h)=0\) in \(F_{2^d}\), then
\(29\mid d\); because \(d\) is even, the extension degree \(d/29\) is even,
so its minimal trace one becomes ambient trace zero.
\end{proof}

The supplement gives the complete exponent lists.  One verifier uses symbolic
Bareiss resultants, Rabin tests, and quotient-field traces.  A code-independent
verifier reconstructs both resultants from all \(256\) evaluations in an AES
model of \(\GF(256)\), using Gaussian Sylvester determinants, interpolation,
Berlekamp nullities, and coefficient traces.  Independent degree bounds
\(240\) and \(214\) are below \(256\); both reconstructed resultants have
degree \(214\).

\begin{corollary}
\label{cor:proper}
An admissible zero fibre satisfying \eqref{eq:zeroexterior} gives a proper
array with \eqref{eq:qparams}.
\end{corollary}
\begin{proof}
The exterior bridge gives Theorem~\ref{thm:twofibre}.  Equation
\eqref{eq:S1}, \eqref{eq:S3}, Lemma~\ref{lem:resultant}, and
\eqref{eq:powerobstruction} force nonconstant multiplicative intersections.
\end{proof}

\section{Fixed curves for the exterior condition}
Put
\[
 a=c=h^2,\qquad t=hr,\qquad U(a)=1+a+a^2,
\]
and define
\begin{align}
K(a,T)={}&1+(1+a)T+a^2U(a)T^3+a^2T^4+a^2T^7\notag\\
&+T^8+U(a)T^9,\label{eq:K}\\
D(a,P)={}&P^8+a^2P^4+1,\label{eq:D}\\
L(a,T,P)={}&P K(a,P)+U(a)T^2D(a,P).\label{eq:L}
\end{align}
For admissible \(a\), the denominator in \eqref{eq:normal} below never
vanishes on \(F\).  In fact
\[
 D(a,p)=(p^4+ap^2+1)^2.
\]
If the factor on the right vanished, then \(a\ne0\) and \(y=p^2/a\) would
satisfy \(y^2+y=a^{-2}\), impossible because
\(\Tr(a^{-2})=\Tr(a^{-1})=1\).
Direct substitution gives
\begin{equation}
 B_a(aw^2)=\frac{pK(a,p)}{D(a,p)},\qquad p=hw.
 \label{eq:normal}
\end{equation}
For admissible \(a\), \(K(a,T)\) has a unique field root \(t\), the nonzero
partner of the zero fibre.

\begin{lemma}
\label{lem:criterion}
The zero fibre is exterior if and only if \(L(a,t,P)\) has no root in \(F\).
\end{lemma}
\begin{proof}
By \eqref{eq:normal}, the forbidden membership \(U(a)t^2\in B_a(F)\) is
equivalent to a field root \(p\) of \(L(a,t,P)\).
\end{proof}

Let \(C\) normalize the projective closure of \(K(a,t)=0\), and let \(X\)
normalize that of \(K(a,t)=L(a,t,p)=0\).

\begin{proposition}
\label{prop:geometry}
The curves \(C\) and \(X\) are geometrically integral over \(\F_2\), and so
are the Artin--Schreier covers
\begin{equation}
 y^2+y=1/a
 \label{eq:AS}
\end{equation}
on both curves.
\end{proposition}
\begin{proof}
The Newton polygon of \(K\), as a polynomial in \(a\) over
\(\F_2((t))\), has coefficient valuations \(0,1,3,3,3\) and one lower
edge from \((0,0)\) to \((4,3)\).  In a factorization, lower Newton polygons
add.  Every nonzero lower edge would therefore have slope \(3/4\), hence
horizontal length divisible by \(4\); the total horizontal length is \(4\).
Thus no two positive-degree factors exist, over any algebraic extension of
the constant field, and \(C\) is geometrically integral.

Put \(t=1+s\).  The lowest homogeneous part at \((a,s)=(0,0)\) is
\(a^2(s+a)\).  The tangent \(s+a=0\) occurs with multiplicity one and gives a
branch on the normalization with \(v(a)=1\).  After base change to
\(\overline{\F}_2\), the function \(1/a\) has an odd pole and cannot be of the
form \(g^2+g\).  Hence the Artin--Schreier cover of \(C\) is geometrically
integral.

At the smooth rational point \(Q=(1,1)\), one has \(U(1)=1\), and the monic
polynomial \(L/U\)
specializes to
\begin{equation}
 P^{10}+P^9+P^5+P+1,
 \label{eq:f10}
\end{equation}
which is irreducible and separable over \(\F_2\).  If \(L/U\) factored over
\(\F_2(C)\), choose monic factors.  Their coefficients are integral over the
DVR \(\mathcal O_{C,Q}\) and lie in its fraction field, hence belong to the
DVR because it is integrally closed.  Reduction at \(Q\) would give a
nontrivial monic factorization of \eqref{eq:f10}.  Hence the fibre product is
integral and degree ten over \(C\).

The irreducible specialization gives a smooth closed point of degree ten.
A second smooth closed point has degree three:
\[
 (a,t,p)=(\alpha,\alpha+1,1),\qquad \alpha^3+\alpha+1=0,
\]
as the two gradients are independent.  For the proper normalization \(X\),
the degree of its full constant field divides every closed-point degree.  It
therefore divides both \(10\) and \(3\), hence equals one.  Since \(\F_2\) is
perfect and the function-field extension is separable, \(X\) is smooth with
exact constant field \(\F_2\), and therefore geometrically integral.

Finally,
\[
 L(0,1,P)=(P+1)^8(P^2+P+1).
\]
At a root \(\omega^2+\omega+1=0\), one has
\(\partial L/\partial P=\omega\neq0\).  It lifts unramified above the branch
with \(v(a)=1\), so \(1/a\) again has an odd pole.  This proves the assertion
for the cover of \(X\).
\end{proof}

\section{Hasse--Weil existence for large even dimension}
We first make the ramification estimate explicit.  For \(Z=C\) or \(X\), let
\(n_Q>0\) be the order of a zero of \(a\) at a geometric point \(Q\).  In the
local Artin--Schreier class, successive replacements by \(g^2+g\) either
remove the pole or reduce it to an odd order \(m_Q\leq n_Q\).  Such a pole has
different exponent \(m_Q+1\leq2n_Q\)
\cite[Corollary~2.3]{anbar-ramification}.  Thus the total different has degree
at most \(2\deg(a:Z\to\mathbf P^1)\), and Riemann--Hurwitz
\cite[Theorem~3.4.13]{stichtenoth} gives
\[
 g(Y_Z)\leq2g(Z)-1+\deg(a:Z\to\mathbf P^1).
\]

The plane degree bound for \(C\) and the preceding inequality give
\begin{equation}
 g(C)\leq45,\qquad g(Y_C)\leq98,
 \label{eq:genusC}
\end{equation}
because \(K\) has total degree at most \(11\) and
\(\deg(a:C\to\mathbf P^1)\leq9\).  The fibre product has projective degree
at most \(11\cdot12=132\).  Castelnuovo's genus bound
\cite[Ch.~IV, Theorem~6.4]{hartshorne}, used here only in the deliberately
weaker plane-form estimate \(g\leq(D-1)(D-2)/2\), gives
\begin{equation}
 g(X)\leq8515,\qquad g(Y_X)\leq17119,
 \label{eq:genusX}
\end{equation}
using \(\deg(a:X\to\mathbf P^1)\leq90\).

Remove \(a=0\), poles of \(a\), and \(U(a)=0\), and denote the resulting open
curves by \(C^\circ\) and \(X^\circ\).  The deleted divisors have total degree
at most \(36\) on \(C\), \(360\) on \(X\), \(72\) on \(Y_C\), and \(720\) on
\(Y_X\).  For either open curve \(Z^\circ\), let \(N_0(Z^\circ)\) count points
with \(\Tr(1/a)=0\).  The equation \(y^2+y=1/a\) has exactly two rational
solutions over such a point and none when the trace is one.  Consequently
\[
 \#Y_Z^\circ(\F_q)=2N_0(Z^\circ),\qquad
 N_1(Z^\circ)=\#Z^\circ(\F_q)-\tfrac12\#Y_Z^\circ(\F_q).
\]

Let \(\widetilde A_q=N_1(C^\circ)\) and
\(\widetilde T_q=N_1(X^\circ)\) be the corresponding counts on the
normalizations.  Applying Hasse--Weil
\cite[Theorem~5.2.3]{stichtenoth} to the four smooth projective curves
gives
\begin{align}
 \widetilde A_q&\geq(q+1)/2-188\sqrt q-36,\label{eq:Aq}\\
 \widetilde T_q&\leq(q+1)/2+34149\sqrt q+360.\label{eq:Tq}
\end{align}
We now account explicitly for the fact that these are normalization-point
counts.  For any set \(S\) of rational
points, put \(m_P=\#\nu^{-1}(P)(\F_q)\), where \(\nu\) is the normalization.
Then \(\#\nu^{-1}(S)(\F_q)=\sum_{P\in S}m_P\).  At a smooth point \(m_P=1\);
at a singular point \(0\leq m_P\leq r_P\), where \(r_P\) is the number of
geometric branches.  Hence
\[
 \#\nu^{-1}(S)(\F_q)-\#S
 \leq\sum_{P\in\operatorname{Sing}}(r_P-1).
\]
For the degree-at-most-\(11\) geometrically integral plane curve, the
normalization exact sequence
\[
 0\longrightarrow\mathcal O_{\overline C}
 \longrightarrow\nu_*\mathcal O_C\longrightarrow\mathcal Q\longrightarrow0
\]
gives
\(\sum_P\delta_P=\operatorname{length}\mathcal Q
=p_a(\overline C)-g(C)\leq45\).
At a point with \(r_P\) geometric branches, reduction of the semilocal
normalization modulo its maximal ideals induces a surjection
\((\nu_*\mathcal O_C/\mathcal O_{\overline C})_P
\twoheadrightarrow\overline{\F}_2^{r_P}/\overline{\F}_2\); hence
\(r_P-1\leq\delta_P\).  Applying the displayed inequality to the admissible
trace-one affine pairs gives
\[
 A_q\geq\widetilde A_q-45.
\]
The same fibre identity in the reverse direction gives
\(\#S-\#\nu^{-1}(S)(\F_q)\leq\#(S\cap\operatorname{Sing})\), because only a
singular rational point can have no rational point above it.  For an
admissible rational triple, direct differentiation gives
\(\partial L/\partial p=D(a,p)\neq0\).  If its base pair is smooth on
\(K=0\), then \(\nabla K=(K_a,K_t,0)\) and
\(\nabla L=(L_a,L_t,D)\) are independent, so the triple is smooth and lifts
uniquely to a rational normalization point.  Hence a rational bad triple
with no rational lift can occur only above a singular base pair.  There are
at most \(45\) such pairs, and the admissible polynomial \(L(a,t,p)\) has
degree ten in \(p\), so allowing all of them separately gives
\[
 T_q\leq\widetilde T_q+10\cdot45=\widetilde T_q+450,
\]
where \(T_q\) is now the number of distinct admissible bad affine triples.

By the exact two-to-one parameterization, the nonvanishing of \(D(a,p)\), and
the bijection \(p=hw\), every bad parameter contributes exactly two
\(p\)-values and every good one contributes none.  Therefore
\begin{equation}
 G_q\geq A_q-T_q/2
 \geq(q+1)/4-17262.5\sqrt q-486.
 \label{eq:Gq}
\end{equation}
This is positive at \(q=2^{34}\) and remains positive thereafter.

\begin{theorem}
\label{thm:largeeven}
For every even \(d\geq34\), a proper array with \eqref{eq:qparams} exists.
\end{theorem}

\section{The finite even basis}
For each remaining even \(d\), represent \(F_{2^d}=\F_2[X]/(m_d)\).
Integers encode binary polynomial coefficients, including the leading bit.

\begin{center}
\begin{tabular}{rrrr}
\toprule
\(d\)&modulus&\(a\)&\(t\)\\
\midrule
4&19&2&9\\
6&97&6&58\\
8&501&15&250\\
10&1627&3&65\\
12&7993&26&3012\\
14&30279&3&6853\\
16&124427&3&49789\\
18&320757&26&159424\\
20&1706235&9&1018552\\
22&5740747&6&1921895\\
24&33394861&6&9911191\\
26&99611673&3&53205065\\
28&336533213&3&177941084\\
30&1179354645&3&252669766\\
32&7185006415&8&164093156\\
\bottomrule
\end{tabular}
\end{center}

For each row, exact polynomial arithmetic verifies that \(m_d\) is
irreducible,
\[
 aU(a)\neq0,\quad \Tr(a^{-1})=1,
\]
\begin{equation}
 \deg\gcd(K(a,T),T^{2^d}-T)=1
 \label{eq:gcdK}
\end{equation}
with unique root the displayed \(t\), and
\begin{equation}
 \gcd(L(a,t,P),P^{2^d}-P)=1.
 \label{eq:gcdL}
\end{equation}
Repeated Frobenius reduction proves these identities without enumerating the
field.  An independent implementation obtains the same degrees.

\begin{theorem}
\label{thm:smalleven}
For every even \(d\), \(4\leq d\leq32\), a proper array with
\eqref{eq:qparams} exists.
\end{theorem}

\section{The all-power theorem}
For \(d=3\), \(t=4\), the following \(7\times8\) array on symbols
\(0,\ldots,13\) is a certificate:
\[
\begin{array}{rrrrrrrr}
0&1&3&12&2&10&11&13\\
1&11&0&2&4&13&12&9\\
11&0&2&8&1&5&13&12\\
3&10&9&0&13&6&4&7\\
6&7&12&10&3&1&5&8\\
9&5&7&4&8&2&6&11\\
8&4&5&6&7&9&3&10
\end{array}
\]
Every symbol occurs four times, with \(CC=3\) and \(RC=4\).  The
row-concurrence multiset is
\[
 \{2^{(6)},4^{(9)},6^{(6)}\}.
\]
The fourteen symbol-incidence vectors occur in seven complementary pairs.
Choosing from each pair the vector absent from the first column and adjoining
the zero vector gives a three-dimensional binary linear space whose seven
nonzero vectors all have weight four.  Thus its column support is the
Sylvester trace design at \(q=8\).

\begin{theorem}[All Sylvester--Hadamard powers]
\label{thm:main}
For every \(k\geq2\), \(t=2^k\), there exists
\[
 \SA(4t-2,t,-,t-1,t:(2t-1)\times2t).
\]
\end{theorem}
\begin{proof}
Set \(d=k+1\), \(q=2^d=2t\).  The displayed array covers \(k=2\).
For even \(k\geq4\), \(d\) is odd and Theorem~\ref{thm:odd} applies.
For odd \(k\geq3\), \(d\) is even and Theorem~\ref{thm:largeeven} or
Theorem~\ref{thm:smalleven} applies.
\end{proof}

\section{Prior-work boundary}
The trace columns \eqref{eq:columns} are the classical Sylvester--Hadamard
design.  With
\[
 A_h=\{c\in F^*:\Tr(ch)=0\},
\]
the paired supports \(A_h\cup\{0\}\) and \(F^*\setminus A_h\) are exactly
those in Gordeev--Öhman \cite[Theorem~5.13]{gordeev-ohman}.  Their result
gives an exact-parameter resolvable unordered triple array but not the proper
row supports or ordering in Theorem~\ref{thm:twofibre}.

Same-parameter triple arrays were also known.  McSorley et al.
\cite{mcsorley} give the \(t=8\) parameters, and Preece et al.
\cite{preece} give the same form when \(2t-1\) is an odd prime power.  Their
row concurrence is constant.  O-polynomials, Subiaco hyperovals, line ovals,
and the two-to-one relation for \(O(x)+ux\) are also prior theory
\cite{flocks-ovals,kolsch,abdukhalikov}.

For the proper object, Jäger et al.\ \cite[Table~2]{jager} record
\((v,r,c)=(30,15,16)\) as \texttt{?,T}: their search and listed constructions
did not produce a proper transposed sesqui array, a triple array was known,
and nonexistence was not proved.  The \(t=8\) instance here addresses that
recorded proper-existence gap; this status statement alone does not establish
global priority.

The recent near-triple-array framework of Gordeev--Markstr\"om--Öhman
\cite{near-triple} includes new constructions and selected enumerations of
proper row--column designs, but its listed proper-design cases do not include
the \(15\times16\), 30-symbol target or the all-power family here.

Thus the qualified novelty claim is only the proper ordered all-power family:
it is not the first Hadamard support, same-parameter unordered design, or
same-parameter triple array.

\section{Reproducibility}
\label{sec:repro}
The supplement provides both resultant verifiers, all symbolic identities,
the geometry checks, fifteen finite-basis certificates, the boundary array,
a single full-chain command, and a SHA-256 manifest.  The independent
resultant coefficient hashes are
\begin{quote}
\ttfamily\small
S3=0: 258eeeeca81b726d3bc13d4ab37edfdc0c118ec5b0e813c7dec46ef9ecf7f241\\
S3=S1\textasciicircum3:
212ed4eed451dbc3f7d078b3a61ffdbf51bceb117f548f3febfa68866d032e59
\end{quote}
The manuscript is internally proof-frozen but externally unreviewed.

\section*{Declaration of AI-assisted research and writing}
This research program used OpenAI Codex extensively for candidate selection,
conjecture formation, symbolic and finite-field derivations, adversarial proof
checking, literature-search assistance, verification-code generation, and
manuscript drafting and revision.  The retained claims are exposed through
human-readable arguments and reproducible artifacts, but AI output is not
treated as an authority and independent expert review remains pending.  Codex
is a tool, not an author.  The named author, Carptopus, accepts responsibility
for the manuscript's accuracy, originality, citations, and integrity.

\section*{License}
\textcopyright\ 2026 Carptopus.  This manuscript is licensed under the
\href{https://creativecommons.org/licenses/by/4.0/}{Creative Commons
Attribution 4.0 International License (CC BY 4.0)}.  The source of record is
the \href{https://github.com/Carptopus/MathExplore-Notes}{MathExplore-Notes
repository}.  Reusers must provide appropriate attribution and indicate
whether changes were made.  No warranties are given.

\begin{thebibliography}{99}
\bibitem{abdukhalikov}
K.~Abdukhalikov, Bent functions and line ovals, arXiv:1609.03138 (2016).
\bibitem{bader-qclan}
L.~Bader, G.~Lunardon, and S.~E. Payne,
On \(q\)-clan geometry, \(q=2^e\),
\emph{Bull. Belg. Math. Soc. Simon Stevin} 1 (1994), 301--328,
\url{https://eudml.org/doc/225475}.
\bibitem{conrad-kloosterman}
K.~Conrad, On Weil's proof of the bound for Kloosterman sums,
\emph{J. Number Theory} 97 (2002), 439--446,
\url{https://doi.org/10.1016/S0022-314X(02)00011-2}.
\bibitem{flocks-ovals}
W.~Cherowitzo, T.~Penttila, I.~Pinneri, and G.~F. Royle,
Flocks and ovals, \emph{Geom. Dedicata} 60 (1996), 17--37,
\url{https://doi.org/10.1007/BF00150865}.
\bibitem{gordeev-ohman}
A.~Gordeev and L.-D. Öhman, Resolvable triple arrays,
\emph{Electron. J. Combin.} 33(2) (2026), P2.35,
\url{https://doi.org/10.37236/14977}.
\bibitem{hartshorne}
R.~Hartshorne, \emph{Algebraic Geometry}, Graduate Texts in Mathematics 52,
Springer, 1977,
\url{https://doi.org/10.1007/978-1-4757-3849-0}.
\bibitem{jager}
G.~Jäger, K.~Markström, D.~Shcherbak, and L.-D. Öhman,
Enumeration and construction of row-column designs,
\emph{J. Combin. Des.} 33 (2025), 357--372,
\url{https://doi.org/10.1002/jcd.21991}.
\bibitem{kolsch}
L.~Kölsch and G.~Kyureghyan,
The classifications of o-monomials and of 2-to-1 binomials are equivalent,
\emph{Des. Codes Cryptogr.} 93 (2025), 961--970,
\url{https://doi.org/10.1007/s10623-024-01463-1}.
\bibitem{mcsorley}
J.~P. McSorley, N.~C.~K. Phillips, W.~D. Wallis, and J.~L. Yucas,
Double arrays, triple arrays and balanced grids,
\emph{Des. Codes Cryptogr.} 35 (2005), 21--45.
\bibitem{near-triple}
A.~Gordeev, K.~Markstr\"om, and L.-D. Öhman,
Near triple arrays, \emph{J. Combin. Theory Ser. A} 219 (2026), 106121,
\url{https://doi.org/10.1016/j.jcta.2025.106121}.
\bibitem{okeefe-thas}
C.~M. O'Keefe and J.~A. Thas,
Collineations of Subiaco and Cherowitzo hyperovals,
\emph{Bull. Belg. Math. Soc. Simon Stevin} 3 (1996), 177--192,
\url{https://projecteuclid.org/euclid.bbms/1105540790}.
\bibitem{preece}
D.~A. Preece, W.~D. Wallis, and J.~L. Yucas,
Paley triple arrays, \emph{Australas. J. Combin.} 33 (2005), 237--246.
\bibitem{xiang}
Q.~Xiang, Recent results on difference sets with classical parameters,
\emph{Finite Fields Appl.} 11 (2005), 622--653.
\bibitem{anbar-ramification}
N.~Anbar, H.~Stichtenoth, and S.~Tutdere,
On ramification in the compositum of function fields,
\emph{Bull. Braz. Math. Soc.} 40 (2009), 539--552,
\url{https://doi.org/10.1007/s00574-009-0026-8}.
\bibitem{stichtenoth}
H.~Stichtenoth, \emph{Algebraic Function Fields and Codes}, 2nd ed.,
Graduate Texts in Mathematics 254, Springer, 2009,
\url{https://doi.org/10.1007/978-3-540-76878-4}.
\end{thebibliography}
\end{document}
